\documentclass[11pt]{article}

\DeclareUnicodeCharacter{0301}{\'e}

\usepackage[margin=1in]{geometry}
\usepackage{multirow}
\usepackage{amssymb,amsmath,amsthm,amsfonts}
\usepackage{bm}
\usepackage{textgreek}
\usepackage{mathtools}
\usepackage{enumitem}
\usepackage[numbers,comma,sort&compress]{natbib}
\usepackage{authblk}
\usepackage{graphicx}
\usepackage[font=small]{caption}
\usepackage{subcaption} % [labelformat=simple]
\usepackage{tablefootnote} 
\usepackage{float}
\usepackage[ruled,vlined,linesnumbered]{algorithm2e}
\usepackage{algorithmic}
\renewcommand{\algorithmicrequire}{\textbf{Input: }}  
\renewcommand{\algorithmicensure}{\textbf{Output: }}
\usepackage{physics}


\usepackage{footnote}
\usepackage{xcolor}
\usepackage{mathrsfs}
\usepackage{bbm}
\usepackage{makecell}
\usepackage{pbox}
\usepackage[colorlinks]{hyperref}
\hypersetup{citecolor=blue}
% \urlstyle{same}
%\captionsetup[figure]{labelfont=bf}

\newtheorem{fact}{Fact}[section] 
\newtheorem{theorem}{Theorem}
\newtheorem{definition}{Definition}
\newtheorem{lemma}{Lemma}
\newtheorem{proposition}{Proposition}
\newtheorem{example}{Example}
\newtheorem{corollary}{Corollary}
\newtheorem{remark}{Remark}
\newtheorem{assumption}{Assumption}
\newcommand{\eqn}[1]{(\ref{eqn:#1})}
\newcommand{\eq}[1]{(\ref{eq:#1})}
\newcommand{\prb}[1]{(\ref{prb:#1})}
\newcommand{\thm}[1]{\hyperref[thm:#1]{Theorem~\ref*{thm:#1}}}
\newcommand{\cor}[1]{\hyperref[cor:#1]{Corollary~\ref*{cor:#1}}}
\newcommand{\defn}[1]{\hyperref[defn:#1]{Definition~\ref*{defn:#1}}}
\newcommand{\lem}[1]{\hyperref[lem:#1]{Lemma~\ref*{lem:#1}}}
\newcommand{\prop}[1]{\hyperref[prop:#1]{Proposition~\ref*{prop:#1}}}
\newcommand{\assum}[1]{\hyperref[assum:#1]{Assumption~\ref*{assum:#1}}}
\newcommand{\fig}[1]{\hyperref[fig:#1]{Figure~\ref*{fig:#1}}}
\newcommand{\tab}[1]{\hyperref[tab:#1]{Table~\ref*{tab:#1}}}
\newcommand{\algo}[1]{\hyperref[algo:#1]{Algorithm~\ref*{algo:#1}}}
\renewcommand{\sec}[1]{\hyperref[sec:#1]{Section~\ref*{sec:#1}}}
\newcommand{\append}[1]{\hyperref[append:#1]{Appendix~\ref*{append:#1}}}
\newcommand{\fac}[1]{\hyperref[fac:#1]{Fact~\ref*{fac:#1}}}
\newcommand{\lin}[1]{\hyperref[lin:#1]{Line~\ref*{lin:#1}}}
\newcommand{\itm}[1]{\ref{itm:#1}}
\newcommand{\bdelta}{\mathbf{\Delta}}
\renewcommand{\arraystretch}{1.25}
\newcommand{\algname}[1]{\textup{\texttt{#1}}}
\renewcommand{\i}{\mathrm{i}}
\renewcommand{\d}{\mathrm{d}}
\renewcommand{\L}{\mathcal{L}}
\newcommand{\ad}{\mathrm{ad}}
\newcommand{\T}{\mathcal{T}}
\newcommand{\B}{\Big}
\renewcommand{\P}{\mathcal{P}}
\newcommand{\simu}{\mathrm{sim}}
\newcommand{\Hc}{H_{\mathrm{control}}}
\newcommand{\He}{H_{\mathrm{eff}}}
\newcommand{\A}{\mathcal{A}}
\renewcommand{\O}{\mathcal{O}}
\newcommand{\tot}{\mathrm{tot}}
\newcommand{\G}{\mathcal{G}}
\newcommand{\E}{\mathcal{E}}
\newcommand{\com}{\mathrm{com}}
\def\>{\rangle}
\def\<{\langle}
\def\trans{^{\top}}

\newcommand{\xz}[1]{{\color{cyan} XZ: #1}}
\newcommand{\tyl}[1]{{\color{blue} Tongyang: #1}}
\newcommand{\zsy}[1]{{\color{teal} Shengyu: #1}}
\newcommand{\sz}[1]{{\color{blue} Shuo: #1}}



\title{Mitigating Physical and Trotter Errors in Hamiltonian Simulation via 1D Extrapolation
}
\author{}
\date{}

\begin{document}

\maketitle
Given an $n$-qubit Hamiltonian $H = \sum_{\gamma=1}^{\Gamma} H_{\gamma}$, we aim to simulate $e^{-\i H T}$ by Trotterization under Markovian physical noise. We consider the following two noise models with a controllable noise strength $\lambda >0$:
\begin{itemize}
    \item After each $e^{-\i H_j \tau}$ gate, \sz{Maybe each ``layer" to be more precise} a quantum channel $e^{\lambda\L_j}$ is applied, where $\L_j$ is a Lindbladian acting the same qubits as $H_j$. \sz{For each layer, both the odd and even $Z_iZ_{i+1}$ cover all the qubits, so do the transverse field $X_i$ layer.}
    \item A continuous noise process described by Lindbladian $\lambda\L$.
\end{itemize}
For simplicity, we consider second-order Trotterization first. Let \[
P_2(t)=e^{-\i H_1 t/2}\cdots e^{-\i H_{\Gamma} t/2} e^{-\i H_{\Gamma} t/2} \cdots e^{-\i H_1 t/2}
\] be the second-order Trotter formula. The simulation is performed by repeatedly applying $P_2(sT)$ for $1/s$ times. 

\textcolor{red}{Yicong: In addition to the number of interpolation nodes, try to prove that the overall sample complexity of ZNE procedure via 1D extrapolation is smaller than that by 2D extrapolation.}

\section{Notations}
For a sequence of operators $\A_1, \ldots, \A_{\Gamma}$, we denote $\prod_{\gamma=1}^\Gamma \A_{\gamma}= A_{\Gamma}\A_{\Gamma-1} \ldots \A_1$ and $\prod_{\gamma=\Gamma}^1 \A_{\gamma}= A_{1}\A_{2} \ldots \A_\Gamma$.
\section{Circuit-Level Noise Model}
We define \begin{align*}
    \P_2(t)=e^{-\i \ad_{H_1} t/2}\cdots e^{-\i \ad_{H_\Gamma} t/2} e^{-\i \ad_{H_\Gamma} t/2} \cdots e^{-\i \ad_{H_1}  t/2}
\end{align*}
to be the adjoint action of $P_2(t)$ and \begin{align*}
    \tilde{\P}_2(t,\lambda)=e^{-\i \ad_{H_1} t/2} e^{\lambda \L_1}\cdots e^{-\i \ad_{H_\Gamma} t/2}e^{\lambda \L_\Gamma} e^{-\i \ad_{H_\Gamma} t/2} e^{\lambda \L_\Gamma}\cdots e^{-\i \ad_{H_1}  t/2}e^{\lambda \L_1}
\end{align*}
be the noisy circuit. \sz{The whole circuit is $\P_2(sT)^{1/s}$}

A much simplified model \sz{Assume noise acts after each Trotter step instead of layer}
\begin{align*}
\tilde{\P}_2(t)=\P_2(t)e^{\lambda \L},
\end{align*}
where $\L$ is a summation of local Lindbldians.

Similar ideas can be found in~\cite{hakkaku2025dataefficienterrormitigationphysical}, where they approximate local noise in a quantum circuit as the global depolarizing noise in large depth regime. \sz{Check whether for each Trotter step or the whole quantum circuit.}
\paragraph{Sparse Pauli-Lindblad Model}
As mentioned above, in real circuit implementation, each $H_\gamma$ covers all the qubits as a layer. The sparse Pauli-Lindblad model fits well with the actual experiments assuming the following form:
\begin{align*}
    e^{\lambda\L} = \exp(\sum_i\lambda_i(P_i\cdot P_i^\dagger- I\cdot I )) =\bigcirc_i\left(w_i \cdot+\left(1-w_i\right) P_i \cdot P_i^{\dagger}\right),
\end{align*}
in which $w_i=\left(1+\mathrm{e}^{-2 \lambda_i}\right) / 2$ and $\bigcirc_{i=1}^n O_i(\cdot)=\left(O_n \circ O_{n-1} \circ \cdots \circ O_1\right)$. Each $P_i$ is only supported on a single qubit or a pair of connected 
qubits, the number of possible channels is polynomial and that is why we call it sparse. Despite its simplicity, this model is able to capture crosstalk errors.

The coefficients $\lambda_i$ varies from layer to layer. In practice, we first learn the error channel for each circuit layer by measuring the fidelity of repeating circuit layers. In small-scale experiments (e.g. 5-qubit or 9-qubit superconducting chip), we can learn more components of $P_i$ considering the actual topology and connectivity since the total number of channels is still accessible.

\subsection{Method 1: Separating Trotter Error and Physical Error}
\sz{We try to bound the difference among $\Tr[O\tilde{\P}_2(sT)\rho_0]$, $\Tr[O\P_2(sT)\rho_0]$ and $\Tr[O\mathcal{V_{\text{ideal}}}\rho_0]$. Then generalize to the long-time simulation with $1/s$ steps. Consider variation-of-parameters formula to give Taylor expansion of the observable averages.}

We remark that noiseless $\P_2(t)$ can be regarded as the initial value $\tilde{\P}_2(t,\lambda=0)$ for Taylor expansion with regard to $\lambda$:
\begin{align*}
    \tilde{\P}_2(t,\lambda) = \P_2(t) + \left.\frac{\dd}{\dd\lambda}\tilde{\P}_2(t,\lambda)\right|_{\lambda=0}\lambda+\cdots+\left.\frac{\dd^n}{\dd\lambda^n}\tilde{\P}_2(t,\lambda)\right|_{\lambda=0}\frac{\lambda^n}{n!} + \int_0^\lambda \dd\lambda_1 \left.\frac{\dd^{n+1}}{\dd\lambda^{n+1}}\tilde{\P}_2(t,\lambda)\right|_{\lambda=\lambda_1}\frac{(\lambda-\lambda_1)^n}{n!},
\end{align*}
where the $n$-th order derivative and its value at $\lambda=0$ take the following form:
\begin{align*}
    \frac{\dd^n}{\dd\lambda^n}\tilde{\P}_2(t,\lambda) = \sum_{n_1+\cdots +n_{2\Gamma}=n}\binom{n}{n_1 \cdots n_{2\Gamma}}e^{-i\ad_{H_1}t/2}e^{\lambda\L}(\L)^{n_1}\cdots e^{-i\ad_{H_1}t/2}e^{\lambda\L}(\L)^{n_{2\Gamma}},\\
    \left.\frac{\dd^n}{\dd\lambda^n}\tilde{\P}_2(t,\lambda)\right|_{\lambda=0} = \sum_{n_1+\cdots +n_{2\Gamma}=n}\binom{n}{n_1 \cdots n_{2\Gamma}}e^{-i\ad_{H_1}t/2}(\L)^{n_1}\cdots e^{-i\ad_{H_1}t/2}(\L)^{n_{2\Gamma}}.
\end{align*}

Similar for $\tilde{\P}_2(sT,\lambda)^{1/s}$, we only need to extend the number of exponential terms in product formula from $2\Gamma$ to $2\Gamma/s$, and substitute $t$ with $sT$. The observable error induced by noise is:
\begin{align*}
    \Tr[O\tilde{\P}_2(sT,\lambda)^{1/s}\rho_0] = \Tr[O\P_2(sT)^{1/s}\rho_0] + \Tr[O\sum_{n_1+\cdots +n_{2\Gamma/s}=1}e^{-i\ad_{H_1}sT/2}(\L)^{n_1}\cdots e^{-i\ad_{H_1}sT/2}(\L)^{n_{2\Gamma/s}} \rho_0]\lambda + \cdots 
\end{align*}

\sz{Issue: the coefficients before $\lambda$ contain $s$, which cannot be extracted into polynomial in effective way.}

On the other hand, former work~\cite{watson2024exponentially} has already analyze the algorithmic error by variation-of-parameters formula:
\begin{align*}
    \Tr[O\P_2(sT)^{1/s}\rho_0] = \Tr[O\mathcal{V}_{\text{ideal}}(T)\rho_0] + C_1 s^p + C_2s^{p+1} + \cdots 
\end{align*}

Combining the two parts together, we have
\begin{align*}
    \Tr[O\tilde{\P}_2(sT,\lambda)^{1/s}\rho_0] = \Tr[O\mathcal{V}_{\text{ideal}}(T)\rho_0] + \operatorname{poly}_1(\lambda, s) +\operatorname{poly}_p(s).
\end{align*}

Then we choose a curve $\lambda=\Lambda(s)$ and well-conditioned interpolation set $\{s_i\}$ to cancel the polynomial terms to higher order. The circuit depth is keen to be reduced from 
\begin{align*}
    2\Gamma/s~\text{with}~1/s= O(\frac{NT^{p+1}}{\epsilon^{p}})\longrightarrow 2\Gamma\max (1/s_i)~\text{with}~\min\{s_i\} = 1/\operatorname{polylog}(1/\epsilon).
\end{align*}
\subsection{Method 2: BCH Formula}
We consider a general $p$-th order product formula of the form
\begin{align*}
    P(t) = \prod_{v=1}^Me^{-\i t a_{v} H_{\gamma_v}},
\end{align*}
where the product is ordered from right to left. Here, $a_{v}$ are real coefficients, $M = \Upsilon\Gamma$ is the total number of unitary gates, and the sequence $\gamma_v$ specifies which Hamiltonian is applied at step $v$. The corresponding ideal channel $\P(t)$ is a superoperator given by
\begin{align*}
    \P(t) = \prod_{v=1}^Me^{-\i t a_{v} \A_{\gamma_v}},
\end{align*}
where $\A_{\gamma}:=\ad_{H_{\gamma}}$ denotes the adjoint action of $H_{\gamma}$.

Suppose that the system is subject to physical noise with a controllable strength $\lambda>0$. Specifically, after each layer of unitary gates, a noise channel $e^{\lambda \L_{v}}$ is applied, where $\L_{v}$ is a Lindbladian. The resulting noisy channel is
\begin{align*}
    \tilde{\P}(t,\lambda)=\prod_{v=1}^M\left( e^{\lambda \L_{v}} e^{-\i t a_{v} \A_{\gamma_v}} \right).
\end{align*}

The Baker-Campbell-Hausdorff (BCH) formula gives a series expansion for the effective generator of $\tilde{\P}(t,\lambda)$:
\begin{align}
\label{eq:bch}
    \tilde{\P}(t,\lambda) = \exp\bigg( \sum_{q=1}^{\infty} \Phi_q(t, \lambda) \bigg),
\end{align}
for sufficiently small $t$ and $\lambda$. The first-order term $\Phi_1$ is the sum of all elementary generators:
\begin{align*}
    \Phi_1(t, \lambda) = \sum_{v=1}^M \left( -\i t a_v \A_{\gamma_v} + \lambda \L_v \right) = -\i t \ad_{H} + \lambda\sum_{v=1}^{M}\L_{v}.
\end{align*} 
The $q$-th order term $\Phi_q$ is defined by \begin{align}
    &\Phi_q(t,\lambda)=\sum_{\substack{p_v,q_v\ge 0\nonumber\\ p_1+q_1+\cdots+p_M+q_M=q }}\bigg(\prod_{v=1}^M \frac{\lambda^{p_v}(-\i ta_v)^{q_v}}{p_v!q_v!} \bigg)\\& \hspace{140pt}\cdot\phi_q(\underbrace{\L_M, \ldots, \L_M}_{p_M}, \underbrace{\A_{\gamma_M}, \cdots, A_{\gamma_M}}_{q_M},  \ldots, \underbrace{\L_1,  \ldots, \L_1}_{p_1}, \underbrace{\A_{\gamma_1},  \ldots, A_{\gamma_1}}_{q_1}) \label{eq:def-Phi-q} \\
    &\phi_q(\mathcal{A}_1, \cdots, \mathcal{A}_q)=\frac{1}{q^2} \sum_{\sigma \in S_q} \frac{(-1)^{d_\sigma}}{\binom{q-1}{d_\sigma}}\left[\mathcal{A}_{\sigma(1)},\left[\mathcal{A}_{\sigma(2)},  \ldots,\left[\mathcal{A}_{\sigma(q-1)}, \mathcal{A}_{\sigma(q)}\right]\right]\right], \nonumber
\end{align}
where $S_q$ denotes the set of permutations of $[q]$ and $d_{\sigma}$ denotes the number of $i\in\{1, \ldots, q-1\}$ such that $\sigma(i)>\sigma(i+1)$. 

The right-hand side of Eq.~\eq{bch} converges if $2t\sum_{\gamma=1}^{\Gamma}\|H_{\gamma}\|+\lambda\sum_{v=1}^{M}\|\L_v\|=\O (1)$, which implies \begin{align}
    \label{eq:bch-converge} 
    t = \O\B(\frac{1}{\sum_{\gamma}\|H_\gamma\|}\B), \quad \lambda = \O\B(\frac{1}{\sum_v\|\L_v\|}\B).
\end{align} 
Then we can apply Cauchy's estimate to upper bound the high-order remainder of the expansion, in a manner similar to the proof of \lem{remainder-f}. This yields an error bound that scales with the $1$-norm of the Hamiltonian $H$. To obtain a commutator-scaling bound, we need to deal with the divergence issue of the BCH formula. If $H$ a local Hamiltonian on a lattice of size $N$, \citet{mizuta2025commutator} showed that truncating the BCH formula to $\O(\log(N/\varepsilon))$ terms provides an $\varepsilon$-approximation of the noiseless product formula for simulation time $t=\O(1/\log(N/\varepsilon))$. This is a regime where the BCH formula is not guaranteed to converge. We extend this result to the noisy case. We will show that if the Hamiltonian $H$ and the noise Lindbldians $\L_v$ are local, the truncation of Eq.~\eq{bch} to $\O(\log(N/\varepsilon))$ terms is an $\varepsilon$-approximation of $\tilde{\P}(t,\lambda)$ if $|t|, \lambda =\O(1/\log(N/\varepsilon))$.

A $k$-local Lindbladian on lattice $\Lambda=[N]$ can be decomposed as  \begin{align*}
    \L=\sum_{X\subseteq \Lambda, |X|\le k } \L_X,
\end{align*}
where $\L_X$ only acts on $X$ nontrivially. The Lindbladian $\L$ is $g$-extensive if for any site $j\in \Lambda$, \begin{align*}
    \sum_{X\subseteq \Lambda, j\in X}\|\L_X\|\le g.
\end{align*} 
\sz{Here the norm of superoperator denotes the limit to stretch the trace norm of density matrix, to some extent diamond norm.} Suppose that \begin{align*}
    \A_\gamma = \sum_{X\subseteq \Lambda} \mathcal{A}_{X}^{\gamma}, \quad \L_v=\sum_{X\subseteq \Lambda} \L^{v}_X
\end{align*}
are $k$-local on a lattice of $N$ qubits. Suppose that $\sum_{\gamma\in[\Gamma]}\A_\gamma=\ad_H$ is $g_h$-extensive and $\sum_{v\in [M]}\L_v=:\L_{\tot}$ is $g_l$-extensive. We further assume that $\A_X^{\gamma}$ commute with each other for fixed $\gamma$, and $\L_{X}^{v}$ commute with each other for fixed $v$. This assumption on $\A_X^{\gamma}$ is common in analysis of Trotter methods subject to physical noise: a noise channel is applied following a layer of commuting gates. The assumption on $\L_X^{v}$ is true for Pauli Lindbladian \begin{align*}
    \L_v\colon\rho\mapsto \sum_{i} \lambda_{v,i} (P_{v,i} \rho P_{v,i}^{\dagger}-\rho),
\end{align*}
where $\lambda_{v,i}\ge 0$ and $P_{v,i}$ are Pauli strings of a weight smaller than $k$. 
\begin{theorem}
\label{thm:bch-approx}
    Given $\varepsilon\in (0,1)$, there exists a truncation order $q_0 = \O(\log(N/\varepsilon))$ such that \begin{align*}
        \Bigg\|\tilde{\P}(t,\lambda)-\exp\bigg( -\i t \ad_H+\lambda\L_{\tot}+\sum_{q=2}^{q_0} \Phi_q(t, \lambda) \bigg)\Bigg\|\le \varepsilon,
    \end{align*}
    for all $t\in \mathbb{R}$, $\lambda\ge0$ such that $q_0 k(\Upsilon g_h|t|+g_l \lambda) = \O(1)$.
\end{theorem}
\begin{proof} 
    \xz{TBD}
\end{proof}
We assume the order of the product formula and the locality parameter $k$ is $\O(1)$. Then, the parameter regime in \thm{bch-approx} becomes \begin{align}
    t= \O\B(\frac{1}{\log(N/\varepsilon) g_h}\B), \quad \lambda = \O\B(\frac{1}{\log(N/\varepsilon) g_l}\B), 
\end{align}
which is not guaranteed to be in the convergence region characterized by Eq.~\eq{bch-converge}. Following \cite{aftab2024multi, mizuta2025commutator}, operator $\Phi_q(t,\lambda)$ can be bounded by \begin{align}
    \|\Phi_q(t,\lambda)\|\le \frac{(q-1)!}{2kq^2} (2k(\Upsilon g_h |t| + g_l \lambda))^{q}N. 
\end{align}

Suppose that the Trotter step size is $sT$ for $s\in (0,1)$ and set the physical noise strength to be $\lambda = c(sT)^{p+1}$ for $c > 0$.  Define the effective generator of one Trotter step after truncation as \begin{align}
    \G(s)&:=-\i  \ad_H+c(sT)^p\L_{\tot}+\frac{1}{sT}\sum_{q=2}^{q_0} \Phi_q(sT, c(sT)^{p+1}) \label{eq:def-G}\\
         &=:-\i \ad_H+\sum_{q=p}^{q_0-1} \E_{q+1} (sT)^q, \label{eq:def-E}
\end{align}
such that \begin{align}
    \|\tilde{\P}(sT,c(sT)^{p+1})-\exp(\G(s) sT)\|\le \varepsilon.
\end{align}
We bound the superoperators $\E_{q}$ in the following lemma. 
\begin{lemma}
    Define a constant
    \begin{align}\label{eq:def-B}
        B = \max\{2k\Upsilon g_h, (2kg_l)^{1/(p+1)}, (cg_l)^{1/(p+1)}/2\}.
    \end{align}
    For any $q\ge p+1$, the superoperator $\E_{q}$ defined in Eq.~\eq{def-E} can be bounded by \begin{align}
        \|\E_q\| \le N(q-1)!(2B)^q.
    \end{align}
\end{lemma}
\begin{proof}
As $\P(t)$ is a $p$-th order product formula, the partial summation in the definition of $\Phi_q(t,\lambda)$ (see Eq.~\eqref{eq:def-Phi-q}) with $p_1, \ldots, p_M=0$ equals $0$. Then, the only $(p+1)$-th order term is $ct^{p+1}\L_{\tot}$, and we have \begin{align*}
    \|\E_{p+1}\| = c\|\L_{\tot}\|\le c\sum_{j\in \Lambda}\sum_{X\subseteq \Lambda: j\in X}\sum_{v\in[M]} \|\L_X^v\| \le c Ng_{l}
\end{align*}
Intuitively, the coefficient operator of any $q$-th order term for $q>p+1$ can be bounded by \begin{align}
    \|\E_q\|\le \sum_{j=0}^{\lfloor q/(p+1)\rfloor} \binom{q-pj}{j} \frac{(q-pj-1)!}{2k(q-pj)^2} (2k\Upsilon g_h)^{q-(p+1)j} (2kg_l)^{j }N.
\end{align}
\xz{TBD}
For any  $q>p+1$, the norm of the coefficient operator $\|\E_q\|$ can be bounded by
\begin{align}
\|\E_q\| &\le \sum_{j=0}^{\lfloor q/(p+1)\rfloor} \binom{q-pj}{j} \frac{(q-pj-1)!}{2k(q-pj)^2} (2k\Upsilon g_h)^{q-(p+1)j} (2kg_l)^{j} N \nonumber \\
&\le \sum_{j=0}^{\lfloor q/(p+1)\rfloor} \binom{q-pj}{j} \frac{(q-pj-1)!}{2k(q-pj)^2} B^q N \nonumber \\
&\le \frac{N(q-1)!}{2k}B^q \sum_{j=0}^{q} \binom{q}{j} \le  N(q-1)!(2B)^q.
\end{align}
For $q=p+1$, we have \begin{align*}
    \|\E_q\| = c\|\L_{\tot}\|\le cNg_l\le   N(2B)^{p+1} \le N(q-1)!(2B)^q.
\end{align*}
Therefore, we can conclude that \begin{align}
     \|\E_q\| \le N(q-1)!(2B)^q
\end{align}
for all $q\ge p+1$.
\end{proof}

With this bound of the coefficients of the error series of $\G(s)$, we can bound the extrapolation error using the variation of parameter formula like Ref.~\cite{aftab2024multi}. 
\begin{lemma}
    Define $\mu_{\com}=2N^{\frac{1}{p+1}}q_0B$. The evolution $e^{\G(s) T}$ admits series expansion in $s$ given by \begin{align*}
        e^{\G(s) T}=-e^{-\i \ad_H T}+ \sum_{q=p}^{\infty}\tilde{\E}_{q} s^q,
    \end{align*}
    where $\tilde{\E}_{q}$ are superoperators bounded by \begin{align*}
       \|\tilde{\E}_{q}\|\le   e(2\mu_{\com}T)^{q(1+1/p)},
    \end{align*}
    if $2\mu_{\com}T\ge 1$, and bounded by \begin{align*}
        \|\tilde{\E}_{q}\|\le e(2\mu_{\com}T)^{q},
    \end{align*}
    if $2\mu_{\com}T\le 1$.
\end{lemma}
\begin{proof}
By Dyson series expansion in the interaction picture, we have 
\begin{align*}
    &e^{\G(s) T}-e^{-\i \ad_H T} \\
    = \ & e^{-\i \ad_H T}  \sum_{j=1}^{\infty} \int_{0}^T \d \tau_1 \int_{0}^{\tau_1}\d \tau_2 \cdots \int_{0}^{\tau_{j-1}}\d \tau_j \prod_{j'=j}^{1}\bigg(e^{\i\ad_H \tau_{j'}}\bigg(\sum_{q=p}^{q_0-1}\E_{q+1}s^qT^q\bigg)e^{-\i \ad_H \tau_{j'}}\bigg) \\
    = \ &  e^{-\i \ad_H T}\sum_{q = p}^{\infty}s^{q} \sum_{j=1}^{\infty} \sum_{\substack{p\le q_1,\ldots,q_j<q_0 \\ q_1+\cdots+q_j=q}} \int_{0}^T \d \tau_1 \int_{0}^{\tau_1}\d \tau_2 \cdots \int_{0}^{\tau_{j-1}}\d \tau_j \prod_{j'=j}^{1}\bigg(e^{\i\ad_H \tau_{j'}}\E_{q_{j'}+1}e^{-\i \ad_H \tau_{j'}}T^{q_{j'}}\bigg)\\
    =:\,& \sum_{q=p}^{\infty}\tilde{\E}_{q} s^q   
\end{align*}
The superoperator $\tilde{\E}_{q}$ can be bounded by \begin{align*}
    \big\|\tilde{\E}_{q}\big\|&\le \sum_{j=1}^{\infty} \sum_{\substack{p\le q_1,\ldots,q_j<q_0 \\ q_1+\cdots+q_j=q}} \int_{0}^T \d \tau_1 \int_{0}^{\tau_1}\d \tau_2 \cdots \int_{0}^{\tau_{j-1}}\d \tau_j \prod_{j'=1}^{j}\big(\big\|\E_{q_{j'}+1}\big\|T^{q_{j'}}\big) \\
    &\le  \sum_{j=1}^{\infty} \sum_{\substack{p\le q_1,\ldots,q_j<q_0 \\ q_1+\cdots+q_j=q}}\frac{T^{q+j}}{j!}\prod_{j'=1}^{j}\big(N(q_{j'}+1)!(2B)^{q_{j'}+1}\big) \\
    &\le \sum_{j=1}^{\infty} \sum_{\substack{p\le q_1,\ldots,q_j<q_0 \\ q_1+\cdots+q_j=q}}\frac{T^{q+j}}{j!}\prod_{j'=1}^{j} \mu_{\com}^{q_{j'}+1} = \sum_{j=1}^{\lfloor q/p\rfloor }\frac{(\mu_{\com}T)^{q+j}}{j!}\B( \sum_{\substack{p\le q_1,\ldots,q_j<q_0 \\ q_1+\cdots+q_j=q}} 1\B),
\end{align*}
where the third line follows from $N(q_{j'}+1)!(2B)^{q_{j'}+1} \le (2N^{\frac{1}{p+1}}q_0B)^{q_{j'}+1}=\mu_{\com}^{q_{j'}+1}$ for any $q_{j'} < q_0$.
The second summation can be bounded by \begin{align*}
    \sum_{\substack{p\le q_1,\ldots,q_j<q_0 \\ q_1+\cdots+q_j=q}} 1 &\le \sum_{\substack{ q_1,\ldots,q_j \ge 0 \\ q_1+\cdots+q_j=q}} 1=\binom{q+j-1}{j-1} \le 2^{q+j}.
\end{align*}
Therefore, for $2\mu_{\com}T\ge 1$, we have \begin{align*}
    \big\|\tilde{\E}_{q}\big\|\le \sum_{j=1}^{\lfloor q/p\rfloor }\frac{(2\mu_{\com}T)^{q+j}}{j!} 
    \le (2\mu_{\com}T)^{q+q/p}\sum_{j=1}^{\infty}\frac{1}{j!}\le e(2\mu_{\com}T)^{q(1+1/p)} ,
\end{align*}
and for $2\mu_{\com}T\le 1$, we have \begin{align*}
    \big\|\tilde{\E}_{q}\big\| \le (2\mu_{\com}T)^{q}\sum_{j=0}^{\infty}\frac{ (2\mu_{\com}T)^{j}}{j!}\le (2\mu_{\com}T)^{q} e^{2\mu_{\com}T}\le e(2\mu_{\com}T)^{q}.
\end{align*}

\end{proof}

The number of Trotter steps in the final algorithm should scales as \begin{align*}
    \tilde{\O}\B((\mu_{\com}T)^{1+\frac{1}{p}}\B) =  \tilde{\O}\B(T^{1+\frac{1}{p}} N^{\frac{1}{p}}\max\B\{g_h^{1+\frac{1}{p}}, g_l^{\frac{1}{p}}\B\}\B).
\end{align*}


\section{Continuous Noise Model}
\textbf{Here we assume that to implement the Trotter step for different step size, we only need to linearly scale the control Hamiltonian of the physical device.} Specifically, each application of $P_2(sT)$ is realized by evolving the quantum computer under a time-dependent control Hamiltonian $\Hc(t/s)$ for $t\in [0, sT_{\simu}]$, where $T_{\simu}$ is the physical time required to execute one Trotter step of simulated time $T$. Formally, we have \begin{align*}
    P_2(sT) = \T \exp(\int_{0}^{sT_{\simu}} - \i \Hc(t/s) \, \d t).
\end{align*}
Concurrently, the physical device is subject to Markovian noise described by a Lindbladian $\lambda \L$. While implementing one Trotter step $P_2(sT)$ of step size $sT$, the dynamics of the physical device's density matrix $\tilde{\rho}_s(t)$ is described by \begin{align*}
    \frac{\d \tilde{\rho}_s(t)}{\d t} = -\i\ad_{\Hc(t/s)}(\tilde{\rho}_s(t))+\lambda \L(\tilde{\rho}_s(t)),
\end{align*}
for $t\in[0, sT_{\simu}]$, where $\ad_H(\rho) :=[H,\rho]$.
We define the ideal time-evolution operator and the time-evolution operator with noise strength $\lambda$ as \begin{align*}
    \Phi_s(t_2,t_1):=\T \exp(\int_{t_1}^{t_2} -\i\ad_{\Hc(t/s)}\,\d t),\quad \tilde{\Phi}_{s,\lambda}(t_2, t_1):=\T \exp(\int_{t_1}^{t_2}(-\i \ad_{\Hc(t/s)}+\lambda \L)\,\d t). 
\end{align*}

\subsection{The Magnus Expansion}
\begin{theorem}[{\cite[Theorem III]{magnus1954}}]
\label{thm:magnus-exp}
    Let $A(t)$ be a function of $t$ in an associative ring, and let \begin{align*}
        Y(t) = \T \exp(\int_{0}^{t} A(s)\,\d s).
    \end{align*}
    Then, if certain unspecified conditions of convergence are satisfied, $Y(t)$ can be written in the form \begin{align*}
        Y(t) = \exp \Omega(t),
    \end{align*}
    where 
    \begin{align*}
        \frac{\d \Omega}{\d t}=\sum_{n=0}^{\infty} \frac{B_n}{n!} \ad_{\Omega}^n A
    \end{align*}
    and $B_n$ are the Bernoulli numbers.
\end{theorem}
Take $A_{s,\lambda}(t) := s(-\i \ad_{\Hc(t)}+\lambda\L)$ in \thm{magnus-exp}, and then we have \begin{align}
    Y_{s,\lambda}(T_{\simu}) & = \T\exp(\int_{0}^{T_\simu} s(-\i \ad_{\Hc(t)}+\lambda\L)\,\d t)\\
    & = \T\exp(\int_{0}^{sT_\simu}(-\i \ad_{\Hc(t)}+\lambda\L)\,\d \tau) \\
    & =\tilde{\Phi}_{s,\lambda}(sT_{\simu}, 0).
\end{align}
By \thm{magnus-exp}, if the convergence conditions are satisfied, $\tilde{\Phi}_{s,\lambda}(sT_{\simu}, 0)$ can be written as \begin{align*}
    \tilde{\Phi}_{s,\lambda}(sT_{\simu}, 0) = \exp \Omega_{s,\lambda}(T_{\simu}).
\end{align*}
Following the proof of \cite[Theorem 6]{blanes2009magnus}, we can show that $\Omega_{s,\lambda}(T_{\simu})$ is an analytic function of $s$ and $\lambda$.
\begin{lemma}
    \label{lem:analytic-phi}
    The function $\varphi(s,\lambda):=\log(Y_{s,\lambda}(T_{\simu}))$ is an analytic function of $s$ and $\lambda$ in the set $B_{\pi}$, with \begin{align*}
        B_{\pi} :=\{s,\lambda \in \mathbb{C}: \int_{0}^{T_{\simu}}\|A_{s,\lambda}(t) \|\,\d t\le \pi \}.
    \end{align*}
\end{lemma}
 
Now we take $\lambda = Cs$. 
\begin{lemma}
\label{lem:remainder-f}
    Given $C> 0$, let $\He(s) := \frac{\varphi(s, Cs)}{sT_{\simu}}$ and \begin{align*}
        f(s):= \exp(\He(s) T_{\simu}).
    \end{align*}
    Let \begin{align*}
        r=\min\Big\{ \frac{1}{4\int_{0}^{T_{\simu}}\|\Hc(t)\|\, \d t}, \frac{1}{2\sqrt{C\|\L\|T_{\simu}}}\Big\}.
    \end{align*}
    For any integer $k\ge 1$, there exists a degree-$k$ polynomial $P_k(s)\colon \mathbb{C}\to \mathbb{C}^{2^n\times 2^n}$ and a remainder $R_k(s)\colon \mathbb{C}\to \mathbb{C}^{2^n\times 2^n}$, such that \begin{align*}
        f(s) = P_k(s)+R_k(s), \text{ and }
        \|R_k(s)\|=O\B(\frac{|s|^{k+1}}{r^k(r-|s|)}\B)
    \end{align*}
    for all $|s|\le \min\{r/2,r^2/(8\pi)\}$. 
\end{lemma}
\begin{proof}
    For $|s|< r$, we have \begin{align}
    \label{eq:converge-radius}
        \int_{0}^{T_{\simu}}\|A_{s,Cs}(t) \|\,\d t\le\int_{0}^{T_{\simu}}(2|s|\|\Hc(t)\|+C|s|\|\L\|)\, \d t\le 1.
    \end{align}
    Let $g(s) = \varphi(s, Cs)$. By \lem{analytic-phi}, $g(s)$ is a holomorphic function of $s$ for $|s|< r$. Since $\lim_{s\to 0}g(s)= 0$, 
    $\He(s)=\frac{g(s)}{sT_{\simu}}$ and $f(s) = \exp(\He(s)T_{\simu})$ are holomorphic functions of $s$ for $|s|< r$. For $s=0$, we define $\He(0) = \lim_{s\to 0} \He(s) = g'(0)/T_{\simu}= -\i HT/T_{\simu}$. \xz{TBD by comparing the series expansion of $\exp(\varphi(s, Cs))$ and noisy $P_2(sT)$.}
    \citet{moan1998efficient} showed that $\|\varphi(s,C s)\|\le \pi$ for  $ \int_{0}^{T_{\simu}}\|A_{s, Cs}(t) \|\,\d t\le 1.08686871$. From Eq.~\eq{converge-radius}, the norm $\|g(s)\|\le \pi $ for all $|s|\le r$. 

    
    In the interaction picture, $f(s)$ can be written as \begin{align*}
        f(s)&=\exp(\He(s) T_{\simu}) \\
        &= \exp(\He(0)T_{\simu})\B(I+\T\exp\B(\int_{0}^{T_{\simu}}e^{-t\He(0)} (\He(s)-\He(0))e^{t\He(0)}\, \d t\B)\B).
    \end{align*}
    Note that $\He(0)=-\i HT/T_{\simu}$ is anti-hermitian, we can bound the norm of $f(s)$ as \begin{align*}
        \|f(s)\|&\le \|\exp(\He(0)T_{\simu})\|\B(1+\B\|\T\exp\B(\int_{0}^{T_{\simu}}e^{-t\He(0)} (\He(s)-\He(0))e^{t\He(0)}\, \d t\B)\B\|\B)\\
        &\le 1+\exp\B(\int_{0}^{T_{\simu}}\B\|e^{-t\He(0)} (\He(s)-\He(0))e^{t\He(0)}\B\|\, \d t\B)\\
        &=1+\exp(\|\He(s)-\He(0)\|T_{\simu}).
    \end{align*}
    To bound the exponent, we first give a bound for $\|g''(s)\|$: 
    \begin{align*}
        \|g''(s)\| = \B\|\frac{1}{2\pi i}\int_{|z|=r} g(z) \frac{2}{(z-s)^3}\, \d z\B\| \le \max_{|z|\le r}\|g(z)\| \frac{2r}{(r-|s|)^3}\le \frac{16\pi}{r^2} ,
    \end{align*}
    for $|s|\le r/2$.
    The exponent $\|\He(s)-\He(0)\|T_{\simu}$ can be bounded as \begin{align*}
        \|(\He(s)-\He(0)\|T_{\simu} &= \frac{1}{|s|}\|g(s)-g'(0)s\|\\
        &=\frac{1}{|s|}\B\|\int_0^s g''(z)(s-z)\, \d z\B\| \\
        &\le \frac{1}{|s|}\max_{|z|\le s}\|g''(z)\| \frac{|s|^2}{2} \\
        &\le \frac{|s|}{2}\frac{16\pi}{r^2} \le 1,
    \end{align*}
    for $|s|\le \min\{r/2,r^2/(8\pi)\}$. Then $\|f(s)\|\le 1+e$ for $s$ in the same region. Then for any integer $k\ge 1$, $f(s)$ can be decomposed as \begin{align*}
        f(s) = P_k(s)+R_k(s),
    \end{align*}
    where $P_k(s)$ is the $k$-th-order Taylor polynomial and $R_k(s)$ is the remainder satisfying \begin{align*}
        \|R(s)\| &=\B\|\frac{s^{k+1}}{2\pi \i}\int_{|z|=r}\frac{f(z)}{z^{k+1}(z-s)}\, \d z\B\|\\
        &\le \max_{|z|\le r}\|f(z)\| \frac{|s|^{k+1}}{r^{k}(r-|s|)}\le (1+e)\frac{|s|^{k+1}}{r^{k}(r-|s|)}.
    \end{align*}
    \xz{Here I directly use some results of $\mathbb{C}\to \mathbb{C}$ holomorphic functions for $\mathbb{C}\to \mathbb{C}^{2^n \times 2^n}$ holomorphic functions, which need to further checked the correctness.}
\end{proof}
The scaling of the Trotter step size in \lem{remainder-f} is similar to qDRIFT, since the first-order noise dominants. Intuitively, if we take the noise strength $\lambda = Cs^{p}$ and use $p$-th-order product formulas, $s=O(r^{1+1/p})$ suffices. The key is to prove that the $g^{(j)}(0)=0$ for $2\le j\le p+1$, where $g(s) = \varphi(s, Cs^p)$.

Take $\lambda = C s^p$ and $\P$ be a $p$-th-order product formula. By the variation of parameter formula, we have \begin{align*}
    Y_{s,Cs^p}(T_{\simu}) & = \T \exp(s\int_0^{T_{\simu}}-\i \ad_{\Hc(t)} \, \d t)\\
    &\quad+\int_{0}^{T_{\simu}} \exp_{\T}(s\int_{t}^{T_{\simu}}-\i \ad_{\Hc(\tau)}\, \d \tau )Cs^{p+1}\L\exp_{\T}(s\int_{0}^{t}(-\i \ad_{\Hc(\tau)}+Cs^p\L)\, \d \tau )\, \d t \\
     &= \P(sT)+O(s^{p+1}) = e^{-\i HsT}+O(s^{p+1})
\end{align*}
Then we have \begin{align*}
    \varphi(s,Cs^{p}) &= \log(Y_{s,Cs^p}(T_{\simu}))\\
    & =\sum_{k=1}^{\infty}\frac{(-1)^{k-1}}{k}(Y_{s,Cs^p}(T_{\simu})-I)^k \\
    & = \log(e^{-\i H s T})+O(s^{p+1})=-\i HsT + O(s^{p+1}).
\end{align*}

\bibliographystyle{plainnat}
\bibliography{ref}
\end{document}
